\pagebreak
\chapter{Live Fast Die Immediately \cite{Hoosain2026} \label{chap:Livefastdie}}
\hrule
\vspace{.5cm}

\section{Abstract and intro}
\hrule
\vspace{.5cm}

\begin{list}{-}{}
\item Seek to better understand BH-accretion disk systems 
\item This includes the spin rate and how might that spin rate be affected by neutrino emission
\end{list}

\section{Collapsers}
\label{Collapsers}
\hrule
\vspace{.5cm}
\begin{list}{-}{}
\item Stars massive enough that there is no \gls{ac:TII-SNe}, instead a direct collapse to a BH
\item This leaves behind the BH with a disk of material which can power strong jets
\item This can prove to be the source of the long duration GrBs 
\end{list}

\section{Neutrino Cooling}
\hrule
\vspace{.5cm}
\begin{list}{-}{}
\item Jets are powered by the spin of the central BH and B field strength
\item For GrBs to be produce the accretion disk needs to be in a \gls{AC:MAD} state. 
\item Neutrinos are thought to be a source of `cooling' or reduction of energy from the systems, as they are produced in the core during collapses and can easily escape, drawing energy away from the system
\item This paper proposes the effects of this neutrino cooling on the system
\item Authors model two versions, one with constant density, and one with a `power slope radius,' meaning that density various with radius. This affects accretion rates and shows that faster spinning systems accrete slower
\item Accretion rates also play a key part in neutrino rate emission efficiency
\item Slower spinning BHs lead to weaker jets, leading to the BH being kicked out of the \gls{AC:MAD} faster, leading to shorter and fainter GrBs
\item Neutrino cooling \textit{doesn't} affect the cooling rate directly, instead effecting other things, such as the magnetic field(????) 
\end{list}